\documentclass[11pt]{article}
\usepackage[a4paper,margin=25mm]{geometry}
\usepackage{amsmath,amssymb,amsthm,booktabs,hyperref,xurl}
\hypersetup{hidelinks}
\setlength{\emergencystretch}{2em}
\newtheorem{proposition}{Proposition}\newtheorem{theorem}[proposition]{Theorem}
\newcommand{\Q}{\mathbb Q}
\title{Two contact endpoints replace two shared power coefficients}
\author{\(\mathbb Z\mathbb Q\) Diophantine Project}
\date{23 November 2026}
\begin{document}
\maketitle

\begin{abstract}
The retained aggregate-spacing identity makes the root block's second
contact the sixteenth point of the common progression. Its centered
square and fourth power can therefore replace two common interpolation
coefficients. A guarded rechart first puts its four contact registers
in the same gauge as the first contact; the recharted inverse and gate
rows remain quadratic. Two affine eliminations then remove two
coordinates, and exact polynomial identities remove two redundant
residuals. The resulting complete general and specified systems have
968 variables with 823 rows and 1143 variables with 994 rows,
respectively. This step is a rational zero-set equivalence and needs
no additional arithmetic completeness or prime-supplier argument.
\end{abstract}

\section{The endpoint already present in the system}

Use the aggregate-spacing systems of \cite{Spacing}, including the
subsequent sign-square deletion on the general route \cite{Branch}.
Write \(\rho\) for the common contact gauge and retain its exact
guards
\[
                 \rho\epsilon=1,\qquad \epsilon_2=\epsilon^2.
\]
The root block is labelled \(b\). Its first and second contact
coordinates are \(t=\texttt{int\_b\_2}\) and
\(t'=\texttt{int\_b\_68}\). The proved spacing is
\[
 D=T_b/16,\qquad d=(t'-t)/16,
\]
where \(T_b\) is the root block's first physical aggregate shift.
In the general chart \(t=\lambda/\rho\),
\(t'=(\lambda+T_b)/\rho\), and \(d=D/\rho\).
In the specified chart the corresponding three coordinates are
unscaled. Thus the already present second contact has physical
argument \(\lambda+16D\). It is not a new contact or a new input
block.

Retain the root centered powers \(u,f\), and denote the common
progression coefficients by \(v,w,a_1,a_2,a_3,a_4\). Their
established formulas are
\begin{equation}\label{interpolation}
 u_i=u+2iv+i^2w,\qquad
 f_i=f+ia_1+i^2a_2+i^3a_3+i^4a_4.
\end{equation}
In the general chart \(u_i=(\lambda+iD-1)^2\) and
\(f_i=(\lambda+iD-1)^4\). In the specified chart
\[
 u_i=\rho\bigl(22(\lambda+iD)-3\bigr)^2,
 \qquad f_i=\rho\bigl(22(\lambda+iD)-3\bigr)^4.
\]
The endpoint's old stored powers use a different gauge. We align
that gauge before applying \eqref{interpolation} at \(i=16\).

\section{A guarded rechart of four private registers}

Reinterpret the root second contact's square, fourth power,
combination and inverse-weight registers. Their literal names are
\[
 \texttt{int\_b\_69},\quad\texttt{int\_b\_70},\quad
 \texttt{int\_b\_71},\quad
 \begin{cases}
 \texttt{int\_b\_158}&\text{general},\\
 \texttt{int\_b\_160}&\text{specified}.
 \end{cases}
\]
For all four coordinates, the recovery of the old value from the new
one is
\begin{equation}\label{rechart}
 \text{old value}=\epsilon^k\,\text{new value},\qquad
 k=\begin{cases}1&\text{general},\\2&\text{specified}.\end{cases}
\end{equation}
The retained unit guard makes this a rational bijection. Denote the
new square and fourth-power values by \(u_*,f_*\).

\begin{proposition}\label{chart}
After \eqref{rechart}, all second-contact rows remain equivalent to
quadratic rows. The new powers \(u_*,f_*\) have exactly the gauge
of the first-contact powers in \eqref{interpolation}.
\end{proposition}
\begin{proof}
The complete occurrence audit finds the four recharted coordinates
only in the following six rows. Substitute \eqref{rechart}, use the
two gauge identities, and multiply by the displayed nonzero factor:
\begin{center}
\begin{tabular}{lc}
\toprule
Second-contact row & Factor multiplying the substituted row\\
\midrule
Centered square & \(1\)\\
Centered fourth power & \(\rho^{2k}\)\\
Contact combination & \(\rho^{2k}\)\\
Contact relation & \(\rho^{2k}\)\\
Inverse weight & \(\rho^k\)\\
Gate & \(1\)\\
\bottomrule
\end{tabular}
\end{center}
The square and fourth-power rows become, respectively,
\[
 \begin{array}{ll}
 \text{general}:&
 t'^2-2\epsilon t'+\epsilon_2-\epsilon_2u_*=0,
 \qquad u_*^2-f_*=0,\\[2pt]
 \text{specified}:&
 (22t'-3)^2-\epsilon u_*=0,
 \qquad u_*^2-\rho f_*=0.
 \end{array}
\]
These are the first-contact chart equations at \(t'\).
The combination and relation rows become literal copies of the
first-contact polynomials with their argument and four private
registers renamed. This statement includes every numerical coefficient;
the independent checker derives it directly by Laurent substitution.

Let \(R_*\) denote the new inverse-weight register, and let \(W_*\)
be the retained second contact weight. The inverse row becomes
\[
 \epsilon_2-W_*R_*=0\quad\text{in the general chart},\qquad
 \epsilon-W_*R_*=0\quad\text{in the specified chart}.
\]
In the general gate row, the old term \(-\epsilon R_{\rm old}\)
becomes \(-\epsilon_2R_*\). In the specified gate row, the old
term \(-\rho R_{\rm old}\) becomes \(-\epsilon R_*\).
Every other gate term is unchanged. These expressions are quadratic.
No native-pole row, point row, contact weight, or gauge guard changes.
This verifies the whole rechart and proves the proposition.
\end{proof}

\section{Two affine eliminations and two exact identities}

At the endpoint \(i=16\), equation \eqref{interpolation} becomes
\[
 u_*=u+32v+256w,\qquad
 f_*=f+16a_1+256a_2+4096a_3+65536a_4.
\]
Eliminate the two common coordinates by the affine substitutions
\begin{equation}\label{eliminate}
 v=\frac{u_*-u-256w}{32},\qquad
 a_1=\frac{f_*-f-256a_2-4096a_3-65536a_4}{16}.
\end{equation}
All occurrences in the entire table receive these substitutions.
Because their coefficients are rational constants, quadratic degree
is preserved.

Here are the exact identities reconstructing the two endpoint power
rows that can now be deleted. Use the unscaled progression residuals
of \cite{Progression}:
\[
 \begin{array}{lll}
 \text{general}:&R_v=\epsilon_2v-d(t-\epsilon),
 &R_w=\epsilon_2w-d^2,\\
 \text{specified}:&R_v=\epsilon v-22d(22t-3),
 &R_w=\epsilon w-484d^2.
 \end{array}
\]
Put \(\chi=1\) in the general chart and \(\chi=\rho\)
in the specified chart, and define
\[
 \begin{aligned}
 R_1&=\chi a_1-4uv,&R_2&=\chi a_2-2uw-4v^2,\\
 R_3&=\chi a_3-4vw,&R_4&=\chi a_4-w^2.
 \end{aligned}
\]
Let \(Q_0,P_0\) be the retained root first-contact square and
fourth-power rows, with \(P_0=u^2-\chi f\).
After \eqref{eliminate} and \(d=(t'-t)/16\), the recharted
endpoint residuals satisfy the literal polynomial identities
\begin{equation}\label{identities}
 \begin{split}
 Q_*&=Q_0-32R_v-256R_w,\\
 P_*&=P_0-16R_1-256R_2-4096R_3-65536R_4,
 \qquad P_*=u_*^2-\chi f_*.
 \end{split}
\end{equation}
The rational constant row normalizations in the stored table are
accounted for by the coefficient certificate. In particular the
existing mixed and spacing rows are \(16R_v,256R_w\) in the
general chart and \(8R_v,64R_w\) in the specified chart.

\begin{theorem}[Exact guarded equivalence]
The rechart \eqref{rechart}, substitutions \eqref{eliminate}, and
deletion of the endpoint square and fourth-power rows give a
bijective rational reconstruction of the predecessor zero set.
All existing source, integer-interface and contact-supplier hypotheses
remain unchanged.
\end{theorem}
\begin{proof}
Given a new solution, assign \(v,a_1\) by \eqref{eliminate}.
Every retained recharted row then holds by its defining substitution.
Equations \eqref{identities} restore the omitted endpoint rows.
Proposition~\ref{chart} recovers all old coordinates by
\eqref{rechart}, using only the already retained nonzero gauge.

Conversely start with a predecessor solution and apply the inverse
of \eqref{rechart}. Its first and second contact arguments differ
by exactly \(16D\), because the aggregate-spacing identity is
already imposed. The retained first power and progression rows
therefore force the endpoint square and fourth power to be the
values in \eqref{interpolation}. Their old defining rows make
\(u_*,f_*\) exactly those values. Hence \eqref{eliminate} holds
for the old \(v,a_1\), and deleting them gives a new solution.
Both directions determine every reconstructed coordinate uniquely.
No arithmetic witness is reselected in this argument.
\end{proof}

\section{The complete tables and their independent audit}

The general predecessor is the frozen sign-square table from
\nolinkurl{undecidability/pell_signed_branch_v101/}; the specified
predecessor is the frozen aggregate-spacing table from
\nolinkurl{undecidability/aggregate_spacing_v100/}. Their exact
censuses are
\begin{center}
\begin{tabular}{lrrrr}
\toprule
Interface & Old variables & New variables & Old rows & New rows\\
\midrule
General & 970 & 968 & 825 & 823\\
Specified & 1145 & 1143 & 996 & 994\\
\bottomrule
\end{tabular}
\end{center}
The ordinary-\(abc\) system is unchanged at 695 variables and
590 quadratic residuals. Each complete system gives the usual
quartic sum of squares. The present equivalence adds no arithmetic
or source-universality hypothesis and makes no new effective claim
about the general fixed parameter.

The deterministic compiler and independent audit are
\begin{verbatim}
python undecidability/contact_endpoint_v102/build_contact_endpoint.py
python undecidability/contact_endpoint_v102/verify_contact_endpoint.py
\end{verbatim}
The compiler copies the first-contact combination templates. The
checker instead reconstructs each old polynomial directly by the
guarded Laurent substitution \(\rho=R,\epsilon=R^{-1}\),
\(\epsilon_2=R^{-2}\) and \eqref{rechart}. It then checks both
affine eliminations and the two omitted-row identities. It audits
all 32 or 28 changed rows, all 791 or 966 unchanged rows, every
occurrence of the four recharted coordinates, both literal gauge
guards, support, degree, integer coefficient rings, metadata, and
hashes. It imports no compiler.

\begin{thebibliography}{9}
\small\setlength{\itemsep}{2pt}
\bibitem{Spacing} \(\mathbb Z\mathbb Q\) Diophantine Project,
\emph{A retained aggregate shift supplies the common contact spacing}.
\par\nolinkurl{undecidability/aggregate_spacing_v100/aggregate_spacing.tex}.
\bibitem{Progression} \(\mathbb Z\mathbb Q\) Diophantine Project,
\emph{A common arithmetic progression of contact arguments and shared
coefficients of their centered powers}.
\par\nolinkurl{undecidability/general_contact_v94/affine_contact_progressions.tex}.
\bibitem{Branch} \(\mathbb Z\mathbb Q\) Diophantine Project,
\emph{Both first-Pell signs have the same coding projection}.
\par\nolinkurl{undecidability/pell_signed_branch_v101/signed_branch_projection.tex}.
\end{thebibliography}
\end{document}
